\section*{Hoofdstuk \ref{ch:g6:earth-a-planet} --- De Aarde, een planeet van het zonnestelsel}

\begin{solution}{exo:g6:earth-a-planet:1}
\underline{Mercurius}, \underline{Venus}, \underline{de Aarde},
\underline{Mars}, Jupiter, Saturnus, Uranus, Neptunus. De \omterm{def:g6:earth-a-planet:planet-types}{rotsplaneten} zijn
de eerste vier.
\end{solution}

\begin{solution}{exo:g6:earth-a-planet:2}
``Een \omterm{def:g6:earth-a-planet:planet-types}{rotsplaneet} is klein en heeft een grote dichtheid, met een vast
oppervlak. Een \omterm{def:g6:earth-a-planet:planet-types}{reuzenplaneet} is heel groot, met een kleine dichtheid, en
bestaat vooral uit gassen en vloeistoffen.''
\end{solution}

\begin{solution}{exo:g6:earth-a-planet:3}
Venus (12\,104~km, tegen 12\,756~km voor de Aarde). De grootste is Jupiter
(142\,984~km); de kleinste is Mercurius (4\,879~km).
\end{solution}

\begin{solution}{exo:g6:earth-a-planet:4}
(a) Niet waar: Venus heeft geen vloeibaar water; het is er veel te heet.
(b) Waar ($-65$~$^\circ$C gemiddeld, tegen 15~$^\circ$C). (c) Niet waar: de
vier planeten die het dichtst bij de Zon staan, zijn de \omterm{def:g6:earth-a-planet:planet-types}{rotsplaneten}; de
\omterm{def:g6:earth-a-planet:planet-types}{reuzenplaneten} zijn de vier verste. (d) Waar: ongeveer 0.7~gram per
kubieke centimeter.
\end{solution}

\begin{solution}{exo:g6:earth-a-planet:5}
$4879 \div 12756 \approx 0.382$, dat is 38~\%: Mercurius is iets meer dan
een derde zo breed als de Aarde.
\end{solution}

\begin{solution}{exo:g6:earth-a-planet:6}
$120536 \div 12756 \approx 9.4$: ongeveer 9 Aardes. Saturnus is een \omterm{def:g6:earth-a-planet:planet-types}{reuzenplaneet}.
\end{solution}

\begin{solution}{exo:g6:earth-a-planet:7}
Planeet~A is een \omterm{def:g6:earth-a-planet:planet-types}{reuzenplaneet}: ze is heel groot en haar dichtheid,
1.1~gram per kubieke centimeter, is klein, zoals die van de \omterm{def:g6:earth-a-planet:planet-types}{reuzenplaneten}
van het zonnestelsel (tussen 0.7 en 1.7). Planeet~B is een \omterm{def:g6:earth-a-planet:planet-types}{rotsplaneet}: ze
is klein en haar dichtheid, 5.0, is groot, zoals die van de \omterm{def:g6:earth-a-planet:planet-types}{rotsplaneten}
(tussen 3.9 en 5.5).
\end{solution}

\begin{solution}{exo:g6:earth-a-planet:8}
(a) $464 - 15 = 449$~graden heter. (b) $15 - (-65) = 80$~graden kouder.
\end{solution}

\begin{solution}{exo:g6:earth-a-planet:9}
Venus heeft een heel dikke atmosfeer van koolstofdioxide, die de warmte
van de Zon vasthoudt als een zware deken en haar niet laat ontsnappen.
Mercurius, dichter bij de Zon, heeft bijna geen lucht, dus kan hij zijn
warmte niet vasthouden.
\end{solution}

\begin{solution}{exo:g6:earth-a-planet:10}
$228 - 150 = 78$~miljoen kilometer verder. $78 \div 150 = 0.52$: de afstand
van Mars is 52~\% groter dan die van de Aarde.
\end{solution}

\begin{solution}{exo:g6:earth-a-planet:11}
Miljarden jaren geleden had Mars \omterm{def:g1:landscapes-around-us:water}{rivieren}, \omterm{def:g1:landscapes-around-us:water}{meren} en \omterm{def:g5:water-shapes-the-land:flood}{overstromingen}: hij was
warmer en natter. Nu is zijn water bevroren, in het ijs van de poolkappen
en onder de grond bij de \omterm{def:g3:continents-and-oceans:poles}{polen}. Vloeibaar water kan nu niet aan zijn
oppervlak blijven, omdat Mars heel koud is ($-65$~$^\circ$C gemiddeld) en
omdat zijn lucht zo ijl is dat water snel bevriest of in waterdamp
verandert.
\end{solution}

\begin{solution}{exo:g6:earth-a-planet:12}
(a) $8.9 \div 9.8 \approx 0.908$, dat is 91~\%. (b) Als een tweelingzus:
bijna even groot (95~\% van de diameter van de Aarde), bijna dezelfde
zwaartekracht, en ook een \omterm{def:g6:earth-a-planet:planet-types}{rotsplaneet}. Niet als een tweelingzus: ze is veel
heter (464~$^\circ$C tegen 15~$^\circ$C), haar lucht is veel dikker en
bestaat uit koolstofdioxide met \omterm{def:g4:clouds-rain-snow:cloud}{wolken} van zuur, en ze heeft geen vloeibaar
water.
\end{solution}

\begin{solution}{pb:g6:earth-a-planet:1}
\textbf{1.} Venus (108~miljoen~km), de Aarde (150), Mars (228).

\textbf{2.} $150 - 108 = 42$: Venus staat 42~miljoen kilometer dichterbij.
$228 - 150 = 78$: Mars staat 78~miljoen kilometer verder.

\textbf{3.} $6792 \div 12756 \approx 0.532$, dat is 53~\%: Mars is ongeveer
half zo breed als de Aarde.

\textbf{4.} Alle drie zijn \omterm{def:g6:earth-a-planet:planet-types}{rotsplaneten}: hun dichtheden (5.2, 5.5 en
3.9~gram per kubieke centimeter) zijn groot, en hun diameters zijn klein
vergeleken met die van de \omterm{def:g6:earth-a-planet:planet-types}{reuzenplaneten}.

\textbf{5.} De dikste lucht is op Venus (92~bar); de ijlste op Mars
(0.01~bar).

\textbf{6.} $1 \div 0.01 = 100$: de lucht van Mars drukt honderd keer minder
dan de lucht van de Aarde.

\textbf{7.} Haar dikke lucht van koolstofdioxide houdt de warmte van de Zon
vast en laat haar niet ontsnappen, als een zware deken.

\textbf{8.} $92 \div 1 = 92$: op de grond van Venus drukt de lucht
\textbf{92 keer} zo hard als op de grond van de Aarde. Een sonde moet heel
sterke wanden hebben om niet verpletterd te worden (en moet ook bestand
zijn tegen de hitte en de zure \omterm{def:g4:clouds-rain-snow:cloud}{wolken}).

\textbf{9.} De gemiddelde temperatuur van de Aarde, 15~$^\circ$C, ligt
tussen 0~$^\circ$C en 100~$^\circ$C: water is niet bevroren en kookt niet,
dus blijft het vloeibaar.

\textbf{10.} Op Venus, bij 464~$^\circ$C, ver boven 100~$^\circ$C, zou het
water meteen wegkoken en in waterdamp veranderen. Op Mars, bij
$-65$~$^\circ$C, ver onder 0~$^\circ$C, zou het tot ijs bevriezen (en
langzaam in waterdamp veranderen, omdat de lucht zo ijl is).

\textbf{11.} Lang geleden was Mars warmer en natter dan nu: er kon vloeibaar
water over zijn oppervlak stromen. Zijn \omterm{def:g5:weather-and-climate:climate}{klimaat} is sterk veranderd: hij is
nu een koude, droge \omterm{def:g5:weather-and-climate:desert}{woestijn}.

\textbf{12.} Bijvoorbeeld: ``Beste sondes, hier op Aarde hebben we
vloeibaar water, in \omterm{def:g2:water-on-earth:ocean}{oceanen}, \omterm{def:g1:landscapes-around-us:water}{rivieren} en \omterm{def:g1:landscapes-around-us:water}{meren}. Onze temperaturen zijn
zacht, gemiddeld 15~$^\circ$C, tussen het bevriezen en het koken van water.
We staan op de juiste afstand van de Zon, niet te dichtbij en niet te ver.
En onze atmosfeer houdt ons warm, geeft levende wezens de gassen die ze
nodig hebben, en laat water vloeibaar blijven. Daarom is hier leven
mogelijk.''
\end{solution}
